\documentclass[]{article}
\usepackage{lmodern}
\usepackage{amssymb,amsmath}
\usepackage{ifxetex,ifluatex}
\usepackage{fixltx2e} % provides \textsubscript
\ifnum 0\ifxetex 1\fi\ifluatex 1\fi=0 % if pdftex
  \usepackage[T1]{fontenc}
  \usepackage[utf8]{inputenc}
\else % if luatex or xelatex
  \ifxetex
    \usepackage{mathspec}
  \else
    \usepackage{fontspec}
  \fi
  \defaultfontfeatures{Ligatures=TeX,Scale=MatchLowercase}
  \newcommand{\euro}{€}
\fi
% use upquote if available, for straight quotes in verbatim environments
\IfFileExists{upquote.sty}{\usepackage{upquote}}{}
% use microtype if available
\IfFileExists{microtype.sty}{%
\usepackage{microtype}
\UseMicrotypeSet[protrusion]{basicmath} % disable protrusion for tt fonts
}{}
\usepackage[margin=1in]{geometry}
\usepackage{hyperref}
\PassOptionsToPackage{usenames,dvipsnames}{color} % color is loaded by hyperref
\hypersetup{unicode=true,
            pdftitle={Partial inverses by Sweeping},
            pdfauthor={Jan de Leeuw},
            pdfborder={0 0 0},
            breaklinks=true}
\urlstyle{same}  % don't use monospace font for urls
\usepackage{color}
\usepackage{fancyvrb}
\newcommand{\VerbBar}{|}
\newcommand{\VERB}{\Verb[commandchars=\\\{\}]}
\DefineVerbatimEnvironment{Highlighting}{Verbatim}{commandchars=\\\{\}}
% Add ',fontsize=\small' for more characters per line
\usepackage{framed}
\definecolor{shadecolor}{RGB}{248,248,248}
\newenvironment{Shaded}{\begin{snugshade}}{\end{snugshade}}
\newcommand{\KeywordTok}[1]{\textcolor[rgb]{0.13,0.29,0.53}{\textbf{{#1}}}}
\newcommand{\DataTypeTok}[1]{\textcolor[rgb]{0.13,0.29,0.53}{{#1}}}
\newcommand{\DecValTok}[1]{\textcolor[rgb]{0.00,0.00,0.81}{{#1}}}
\newcommand{\BaseNTok}[1]{\textcolor[rgb]{0.00,0.00,0.81}{{#1}}}
\newcommand{\FloatTok}[1]{\textcolor[rgb]{0.00,0.00,0.81}{{#1}}}
\newcommand{\ConstantTok}[1]{\textcolor[rgb]{0.00,0.00,0.00}{{#1}}}
\newcommand{\CharTok}[1]{\textcolor[rgb]{0.31,0.60,0.02}{{#1}}}
\newcommand{\SpecialCharTok}[1]{\textcolor[rgb]{0.00,0.00,0.00}{{#1}}}
\newcommand{\StringTok}[1]{\textcolor[rgb]{0.31,0.60,0.02}{{#1}}}
\newcommand{\VerbatimStringTok}[1]{\textcolor[rgb]{0.31,0.60,0.02}{{#1}}}
\newcommand{\SpecialStringTok}[1]{\textcolor[rgb]{0.31,0.60,0.02}{{#1}}}
\newcommand{\ImportTok}[1]{{#1}}
\newcommand{\CommentTok}[1]{\textcolor[rgb]{0.56,0.35,0.01}{\textit{{#1}}}}
\newcommand{\DocumentationTok}[1]{\textcolor[rgb]{0.56,0.35,0.01}{\textbf{\textit{{#1}}}}}
\newcommand{\AnnotationTok}[1]{\textcolor[rgb]{0.56,0.35,0.01}{\textbf{\textit{{#1}}}}}
\newcommand{\CommentVarTok}[1]{\textcolor[rgb]{0.56,0.35,0.01}{\textbf{\textit{{#1}}}}}
\newcommand{\OtherTok}[1]{\textcolor[rgb]{0.56,0.35,0.01}{{#1}}}
\newcommand{\FunctionTok}[1]{\textcolor[rgb]{0.00,0.00,0.00}{{#1}}}
\newcommand{\VariableTok}[1]{\textcolor[rgb]{0.00,0.00,0.00}{{#1}}}
\newcommand{\ControlFlowTok}[1]{\textcolor[rgb]{0.13,0.29,0.53}{\textbf{{#1}}}}
\newcommand{\OperatorTok}[1]{\textcolor[rgb]{0.81,0.36,0.00}{\textbf{{#1}}}}
\newcommand{\BuiltInTok}[1]{{#1}}
\newcommand{\ExtensionTok}[1]{{#1}}
\newcommand{\PreprocessorTok}[1]{\textcolor[rgb]{0.56,0.35,0.01}{\textit{{#1}}}}
\newcommand{\AttributeTok}[1]{\textcolor[rgb]{0.77,0.63,0.00}{{#1}}}
\newcommand{\RegionMarkerTok}[1]{{#1}}
\newcommand{\InformationTok}[1]{\textcolor[rgb]{0.56,0.35,0.01}{\textbf{\textit{{#1}}}}}
\newcommand{\WarningTok}[1]{\textcolor[rgb]{0.56,0.35,0.01}{\textbf{\textit{{#1}}}}}
\newcommand{\AlertTok}[1]{\textcolor[rgb]{0.94,0.16,0.16}{{#1}}}
\newcommand{\ErrorTok}[1]{\textcolor[rgb]{0.64,0.00,0.00}{\textbf{{#1}}}}
\newcommand{\NormalTok}[1]{{#1}}
\usepackage{graphicx,grffile}
\makeatletter
\def\maxwidth{\ifdim\Gin@nat@width>\linewidth\linewidth\else\Gin@nat@width\fi}
\def\maxheight{\ifdim\Gin@nat@height>\textheight\textheight\else\Gin@nat@height\fi}
\makeatother
% Scale images if necessary, so that they will not overflow the page
% margins by default, and it is still possible to overwrite the defaults
% using explicit options in \includegraphics[width, height, ...]{}
\setkeys{Gin}{width=\maxwidth,height=\maxheight,keepaspectratio}
\setlength{\parindent}{0pt}
\setlength{\parskip}{6pt plus 2pt minus 1pt}
\setlength{\emergencystretch}{3em}  % prevent overfull lines
\providecommand{\tightlist}{%
  \setlength{\itemsep}{0pt}\setlength{\parskip}{0pt}}
\setcounter{secnumdepth}{5}

%%% Use protect on footnotes to avoid problems with footnotes in titles
\let\rmarkdownfootnote\footnote%
\def\footnote{\protect\rmarkdownfootnote}

%%% Change title format to be more compact
\usepackage{titling}

% Create subtitle command for use in maketitle
\newcommand{\subtitle}[1]{
  \posttitle{
    \begin{center}\large#1\end{center}
    }
}

\setlength{\droptitle}{-2em}
  \title{Partial inverses by Sweeping}
  \pretitle{\vspace{\droptitle}\centering\huge}
  \posttitle{\par}
  \author{Jan de Leeuw}
  \preauthor{\centering\large\emph}
  \postauthor{\par}
  \predate{\centering\large\emph}
  \postdate{\par}
  \date{March 21, 2016}



% Redefines (sub)paragraphs to behave more like sections
\ifx\paragraph\undefined\else
\let\oldparagraph\paragraph
\renewcommand{\paragraph}[1]{\oldparagraph{#1}\mbox{}}
\fi
\ifx\subparagraph\undefined\else
\let\oldsubparagraph\subparagraph
\renewcommand{\subparagraph}[1]{\oldsubparagraph{#1}\mbox{}}
\fi

\begin{document}
\maketitle

{
\setcounter{tocdepth}{2}
\tableofcontents
}
\section{Simple Sweep}\label{simple-sweep}

Suppose \(A\) is a square non-singular matrix. A \emph{simple sweep} on
the \(k^{th}\) element replaces \(A\) by the matrix with elements

\begin{align*}
a_{kk}&=\frac{1}{a_{kk}},\\
a_{ki}&=-\frac{a_{ki}}{a_{kk}},\\
a_{ik}&=\frac{a_{ik}}{a_{kk}},\\
a_{ij}&=a_{ij}-\frac{a_{ik}a_{kj}}{a_{kk}}
\end{align*}

for \(j\not= k\) and \(i\not= k\).

In R that becomes

\begin{Shaded}
\begin{Highlighting}[]
\NormalTok{sweep <-}\StringTok{ }\NormalTok{function (a, k) \{}
  \NormalTok{p <-}\StringTok{ }\NormalTok{a [k, k] }
  \NormalTok{a [k, k] <-}\StringTok{ }\DecValTok{1} \NormalTok{/p}
  \NormalTok{a [-k, -k] <-}\StringTok{ }\NormalTok{a [-k, -k] -}\StringTok{ }\KeywordTok{outer} \NormalTok{(a [-k, k], a [k, -k]) /}\StringTok{ }\NormalTok{p}
  \NormalTok{a [k, -k] <-}\StringTok{ }\NormalTok{-a[k, -k] /}\StringTok{ }\NormalTok{p}
  \NormalTok{a [-k, k] <-}\StringTok{ }\NormalTok{a[-k, k] /}\StringTok{ }\NormalTok{p}
  \KeywordTok{return} \NormalTok{(a)}
\NormalTok{\}}
\end{Highlighting}
\end{Shaded}

\section{Example Sweep}\label{example-sweep}

\begin{Shaded}
\begin{Highlighting}[]
\KeywordTok{set.seed}\NormalTok{(}\DecValTok{12345}\NormalTok{)}
\KeywordTok{mprint} \NormalTok{(a <-}\StringTok{ }\KeywordTok{matrix} \NormalTok{(}\KeywordTok{sample}\NormalTok{(}\DecValTok{1}\NormalTok{:}\DecValTok{9}\NormalTok{), }\DecValTok{3}\NormalTok{, }\DecValTok{3}\NormalTok{))}
\end{Highlighting}
\end{Shaded}

\begin{verbatim}
##      [,1]  [,2]  [,3] 
## [1,]  7.00  9.00  4.00
## [2,]  8.00  3.00  2.00
## [3,]  6.00  1.00  5.00
\end{verbatim}

\begin{Shaded}
\begin{Highlighting}[]
\KeywordTok{mprint} \NormalTok{(}\KeywordTok{sweep}\NormalTok{(a, }\DecValTok{1}\NormalTok{))}
\end{Highlighting}
\end{Shaded}

\begin{verbatim}
##      [,1]  [,2]  [,3] 
## [1,]  0.14 -1.29 -0.57
## [2,]  1.14 -7.29 -2.57
## [3,]  0.86 -6.71  1.57
\end{verbatim}

\subsection{Using Sweeping for the
Inverse}\label{using-sweeping-for-the-inverse}

\begin{Shaded}
\begin{Highlighting}[]
\NormalTok{sweepSubset <-}\StringTok{ }\NormalTok{function (a, k, }\DataTypeTok{verbose =} \OtherTok{FALSE}\NormalTok{) \{}
  \NormalTok{for (i in k) \{}
    \NormalTok{a <-}\StringTok{ }\KeywordTok{sweep} \NormalTok{(a, i)}
    \NormalTok{if (verbose) }\KeywordTok{mprint} \NormalTok{(a)}
  \NormalTok{\}}
  \KeywordTok{return} \NormalTok{(a)}
\NormalTok{\}}
\end{Highlighting}
\end{Shaded}

\begin{Shaded}
\begin{Highlighting}[]
\KeywordTok{mprint} \NormalTok{(b <-}\StringTok{ }\KeywordTok{cSweep} \NormalTok{(a, }\DecValTok{1} \NormalTok{:}\StringTok{ }\KeywordTok{nrow} \NormalTok{(a)))}
\end{Highlighting}
\end{Shaded}

\begin{verbatim}
##      [,1]  [,2]  [,3] 
## [1,]  0.05  0.16  0.03
## [2,]  0.16  0.08 -0.09
## [3,]  0.03 -0.09 -0.26
\end{verbatim}

\begin{Shaded}
\begin{Highlighting}[]
\KeywordTok{mprint} \NormalTok{(a %*%}\StringTok{ }\NormalTok{b)}
\end{Highlighting}
\end{Shaded}

\begin{verbatim}
##      [,1]  [,2]  [,3] 
## [1,]  1.87  1.47 -1.71
## [2,]  0.91  1.33 -0.59
## [3,]  0.57  0.56 -1.23
\end{verbatim}

\subsection{Using Sweeping for
Regression}\label{using-sweeping-for-regression}

\begin{Shaded}
\begin{Highlighting}[]
\NormalTok{x <-}\StringTok{ }\KeywordTok{matrix} \NormalTok{(}\KeywordTok{rnorm} \NormalTok{(}\DecValTok{20}\NormalTok{), }\DecValTok{10}\NormalTok{, }\DecValTok{2}\NormalTok{)}
\NormalTok{y <-}\StringTok{ }\KeywordTok{rnorm} \NormalTok{(}\DecValTok{10}\NormalTok{)}
\KeywordTok{mprint} \NormalTok{(b<-}\KeywordTok{qr.solve} \NormalTok{(x, y))}
\end{Highlighting}
\end{Shaded}

\begin{verbatim}
## [1]  0.11 -0.07
\end{verbatim}

\begin{Shaded}
\begin{Highlighting}[]
\KeywordTok{mprint} \NormalTok{(}\KeywordTok{sum}\NormalTok{((y-x%*%b) ^}\StringTok{ }\DecValTok{2}\NormalTok{))}
\end{Highlighting}
\end{Shaded}

\begin{verbatim}
## [1]  6.39
\end{verbatim}

\begin{Shaded}
\begin{Highlighting}[]
\NormalTok{h <-}\StringTok{ }\KeywordTok{crossprod} \NormalTok{(}\KeywordTok{cbind} \NormalTok{(x,y))}
\KeywordTok{mprint} \NormalTok{(}\KeywordTok{sweepSubset} \NormalTok{(h, }\DecValTok{1} \NormalTok{:}\StringTok{ }\DecValTok{2}\NormalTok{, }\DataTypeTok{verbose =} \OtherTok{TRUE}\NormalTok{))}
\end{Highlighting}
\end{Shaded}

\begin{verbatim}
##               y    
##    0.05  0.03 -0.11
##   -0.03 13.31 -0.95
## y  0.11 -0.95  6.46
##               y    
##    0.05  0.00 -0.11
##    0.00  0.08  0.07
## y  0.11 -0.07  6.39
##               y    
##    0.05  0.00 -0.11
##    0.00  0.08  0.07
## y  0.11 -0.07  6.39
\end{verbatim}

\subsection{The partial inverse}\label{the-partial-inverse}

Suppose \(k\) is a subset of \(\{1,2,\cdots,n\}\) and \(\overline{k}\)
is the complimentary subset. In our examples we usually take, for
convenience, \(k=\{1,2,\cdots,k\}\). The \emph{partial inverse} of \(A\)
on \(k\), which generalizes the simple sweep, is \[
\mathbf{inv}_k(A)=\begin{bmatrix}
A_{kk}^{-1}&-A_{kk}^{-1}A_{k\overline{k}}\\
A_{\overline k k}A_{kk}^{-1}&A_{\overline k\overline k}-A_{\overline k k}A_{kk}^{-1}A_{k\overline{k}}
\end{bmatrix}.
\]

Why is this called a partial inverse ? If \[
\begin{bmatrix}
A_{kk}&A_{k\overline k}\\
A_{\overline k k}&A_{\overline k\overline k}
\end{bmatrix}
\begin{bmatrix}
x\\y
\end{bmatrix}=
\begin{bmatrix}
u\\v
\end{bmatrix},
\] then \[
\begin{bmatrix}
A_{kk}^{-1}&-A_{kk}^{-1}A_{k\overline{k}}\\
A_{\overline k k}A_{kk}^{-1}&A_{\overline k\overline k}-A_{\overline k k}A_{kk}^{-1}A_{k\overline{k}}
\end{bmatrix}
\begin{bmatrix}
u\\y
\end{bmatrix}=
\begin{bmatrix}
x\\v
\end{bmatrix},
\] And conversely. Of course all this assumes that \(A_{kk}\) is
non-singular. The following results are true for an arbitrary index set
\(k\) and its complement \(\overline k\).

\begin{enumerate}
\def\labelenumi{\arabic{enumi}.}
\tightlist
\item
  \(A=\mathbf{inv}_k(\mathbf{inv}_k(A))\)
\item
  \(\mathbf{inv}_k(A)=(\mathbf{inv}_{\overline k}(A))^{-1}\)
\item
  \(\mathbf{inv}_{\overline k}(\mathbf{inv}_k(A))=A^{-1}\)
\item
  \(\mathbf{inv}_{\ell}(\mathbf{inv}_k(A))=\mathbf{inv}_{k}(\mathbf{inv}_\ell(A))\)
\item
  \(\mathbf{inv}_{k}(A)=\mathbf{inv}_{\overline k}(A^{-1})\)
\item
  \((\mathbf{inv}_{k}(A))^{-1}=\mathbf{inv}_{k}(A^{-1})\)
\end{enumerate}

\section{Timing}\label{timing}

\begin{Shaded}
\begin{Highlighting}[]
\KeywordTok{set.seed}\NormalTok{(}\DecValTok{12345}\NormalTok{)}
\NormalTok{b <-}\StringTok{ }\KeywordTok{crossprod}\NormalTok{(}\KeywordTok{matrix}\NormalTok{(}\KeywordTok{rnorm}\NormalTok{(}\DecValTok{1000000}\NormalTok{),}\DecValTok{1000}\NormalTok{,}\DecValTok{1000}\NormalTok{))/}\DecValTok{1000}
\end{Highlighting}
\end{Shaded}

\begin{verbatim}
## [1] "element sweeping in pure R"
##    user  system elapsed 
##  14.974   1.510  15.770 
##    user  system elapsed 
##  14.399   1.377  14.976 
##    user  system elapsed 
##  14.334   1.359  14.881 
##    user  system elapsed 
##  14.388   1.360  14.951 
##    user  system elapsed 
##  14.362   1.363  14.928 
## [1] "element sweeping in R and fortran from R"
##    user  system elapsed 
##   1.423   0.310   1.736 
##    user  system elapsed 
##   1.378   0.301   1.684 
##    user  system elapsed 
##   1.244   0.348   1.595 
##    user  system elapsed 
##   1.309   0.262   1.575 
##    user  system elapsed 
##   1.263   0.267   1.534 
## [1] "element sweeping in R and C from R and fortran from C"
##    user  system elapsed 
##   0.453   0.011   0.468 
##    user  system elapsed 
##   0.420   0.009   0.431 
##    user  system elapsed 
##   0.416   0.011   0.429 
##    user  system elapsed 
##   0.426   0.011   0.439 
##    user  system elapsed 
##   0.435   0.011   0.449 
## [1] "just solve it"
##    user  system elapsed 
##   0.211   0.006   0.071 
##    user  system elapsed 
##   0.197   0.004   0.071 
##    user  system elapsed 
##   0.205   0.003   0.069 
##    user  system elapsed 
##   0.206   0.003   0.070 
##    user  system elapsed 
##   0.197   0.003   0.070
\end{verbatim}

\section{Code}\label{code}

\subsection{R Code}\label{r-code}

\begin{Shaded}
\begin{Highlighting}[]
\NormalTok{rSweep <-}\StringTok{ }\NormalTok{function (a, indi) \{}
\NormalTok{for (j in indi) \{}
    \NormalTok{pv <-}\StringTok{ }\NormalTok{a[j, j]}
    \NormalTok{if (pv ==}\StringTok{ }\DecValTok{0}\NormalTok{) next ()}
    \NormalTok{pr <-}\StringTok{ }\NormalTok{a[j, -j]}
    \NormalTok{pc <-}\StringTok{ }\NormalTok{a[-j, j]}
    \NormalTok{a[j, j] <-}\StringTok{ }\NormalTok{-}\DecValTok{1} \NormalTok{/}\StringTok{ }\NormalTok{pv}
    \NormalTok{a[j, -j] <-}\StringTok{ }\NormalTok{pr /}\StringTok{ }\NormalTok{pv}
    \NormalTok{a[-j, j] <-}\StringTok{ }\NormalTok{pc /}\StringTok{ }\NormalTok{pv}
    \NormalTok{a[-j, -j] <-}\StringTok{ }\NormalTok{a[-j, -j] -}\StringTok{ }\KeywordTok{outer}\NormalTok{(pc, pr) /}\StringTok{ }\NormalTok{pv}
    \NormalTok{\}}
\KeywordTok{return} \NormalTok{(a)}
\NormalTok{\}}

\NormalTok{fSweep <-}\StringTok{ }\NormalTok{function (b , ind, }\DataTypeTok{eps =} \FloatTok{1e-10}\NormalTok{, }\DataTypeTok{trianin =} \OtherTok{FALSE}\NormalTok{, }\DataTypeTok{trianout =} \OtherTok{FALSE}\NormalTok{) \{}
    \NormalTok{if (trianin) \{}
        \NormalTok{m <-}\StringTok{ }\KeywordTok{length} \NormalTok{(b)}
        \NormalTok{n <-}\StringTok{ }\KeywordTok{as.integer} \NormalTok{( (}\KeywordTok{sqrt} \NormalTok{(}\DecValTok{1} \NormalTok{+}\StringTok{ }\DecValTok{8} \NormalTok{*}\StringTok{ }\NormalTok{m) -}\StringTok{ }\DecValTok{1}\NormalTok{) /}\StringTok{ }\DecValTok{2}\NormalTok{)}
        \NormalTok{a <-}\StringTok{ }\NormalTok{b}
    \NormalTok{\} else \{}
        \NormalTok{n <-}\StringTok{ }\KeywordTok{nrow} \NormalTok{(b)}
        \NormalTok{m <-}\StringTok{ }\NormalTok{(n *}\StringTok{ }\NormalTok{n +}\StringTok{ }\NormalTok{n) /}\StringTok{ }\DecValTok{2}
        \NormalTok{a <-}\StringTok{ }\NormalTok{b[}\KeywordTok{upper.tri} \NormalTok{(b, }\DataTypeTok{diag =} \OtherTok{TRUE}\NormalTok{)]}
    \NormalTok{\}}
    \NormalTok{for (k in ind) \{}
        \NormalTok{a <-}\StringTok{ }\KeywordTok{.Fortran}\NormalTok{(}\StringTok{"gsweep"}\NormalTok{,}
            \DataTypeTok{diago =} \KeywordTok{as.double} \NormalTok{(a[}\KeywordTok{cumsum} \NormalTok{(}\DecValTok{1}\NormalTok{:n)]),}
            \DataTypeTok{trian =} \KeywordTok{as.double} \NormalTok{(a),}
            \DataTypeTok{k =} \KeywordTok{as.integer} \NormalTok{(k),}
            \DataTypeTok{l =} \KeywordTok{as.integer} \NormalTok{(}\DecValTok{0}\NormalTok{),}
            \DataTypeTok{m =} \KeywordTok{as.integer} \NormalTok{(m),}
            \DataTypeTok{n =} \KeywordTok{as.integer} \NormalTok{(n),}
            \DataTypeTok{e =} \KeywordTok{as.double} \NormalTok{(eps),}
            \DataTypeTok{ifault =} \KeywordTok{as.integer} \NormalTok{(}\DecValTok{0}\NormalTok{))$trian}
        \NormalTok{\}}
    \NormalTok{if (trianout) \{}
        \KeywordTok{return} \NormalTok{(a)}
    \NormalTok{\} else \{}
        \KeywordTok{return} \NormalTok{(}\KeywordTok{tri2sym} \NormalTok{(a))}
    \NormalTok{\}}
\NormalTok{\}}

\NormalTok{cSweep <-}\StringTok{ }\NormalTok{function (b , ind, }\DataTypeTok{eps =} \FloatTok{1e-10}\NormalTok{, }\DataTypeTok{trianin =} \OtherTok{FALSE}\NormalTok{, }\DataTypeTok{trianout =} \OtherTok{FALSE}\NormalTok{) \{}
    \NormalTok{if (trianin) \{}
        \NormalTok{m <-}\StringTok{ }\KeywordTok{length} \NormalTok{(b)}
        \NormalTok{n <-}\StringTok{ }\KeywordTok{as.integer} \NormalTok{( (}\KeywordTok{sqrt} \NormalTok{(}\DecValTok{1} \NormalTok{+}\StringTok{ }\DecValTok{8} \NormalTok{*}\StringTok{ }\NormalTok{m) -}\StringTok{ }\DecValTok{1}\NormalTok{) /}\StringTok{ }\DecValTok{2}\NormalTok{)}
        \NormalTok{a <-}\StringTok{ }\NormalTok{b}
    \NormalTok{\} else \{}
        \NormalTok{n <-}\StringTok{ }\KeywordTok{nrow} \NormalTok{(b)}
        \NormalTok{m <-}\StringTok{ }\NormalTok{(n *}\StringTok{ }\NormalTok{n +}\StringTok{ }\NormalTok{n) /}\StringTok{ }\DecValTok{2}
        \NormalTok{a <-}\StringTok{ }\NormalTok{b[}\KeywordTok{upper.tri} \NormalTok{(b, }\DataTypeTok{diag =} \OtherTok{TRUE}\NormalTok{)]}
    \NormalTok{\}}
    \NormalTok{a <-}\StringTok{ }\KeywordTok{.C}\NormalTok{(}\StringTok{"hsweep"}\NormalTok{,}
        \DataTypeTok{diago =} \KeywordTok{as.double} \NormalTok{(a[}\KeywordTok{cumsum} \NormalTok{(}\DecValTok{1}\NormalTok{:n)]),}
        \DataTypeTok{trian =} \KeywordTok{as.double} \NormalTok{(a),}
        \DataTypeTok{n =} \KeywordTok{as.integer} \NormalTok{(n),}
        \DataTypeTok{m =} \KeywordTok{as.integer} \NormalTok{(m),}
        \DataTypeTok{ind =} \KeywordTok{as.integer} \NormalTok{(ind),}
        \DataTypeTok{npiv =} \KeywordTok{as.integer} \NormalTok{(}\KeywordTok{length} \NormalTok{(ind)),}
        \DataTypeTok{eps =} \KeywordTok{as.double} \NormalTok{(eps))$trian}
    \NormalTok{if (trianout) \{}
        \KeywordTok{return} \NormalTok{(a)}
    \NormalTok{\} else \{}
        \KeywordTok{return} \NormalTok{(}\KeywordTok{tri2sym} \NormalTok{(a))}
    \NormalTok{\}}
\NormalTok{\}}

\NormalTok{tri2sym <-}\StringTok{ }\NormalTok{function (a) \{}
    \NormalTok{m <-}\StringTok{ }\KeywordTok{length} \NormalTok{(a)}
    \NormalTok{n <-}\StringTok{ }\KeywordTok{as.integer} \NormalTok{( (}\KeywordTok{sqrt} \NormalTok{(}\DecValTok{1} \NormalTok{+}\StringTok{ }\DecValTok{8} \NormalTok{*}\StringTok{ }\NormalTok{m) -}\StringTok{ }\DecValTok{1}\NormalTok{) /}\StringTok{ }\DecValTok{2}\NormalTok{)}
    \NormalTok{b <-}\StringTok{ }\KeywordTok{matrix} \NormalTok{(}\DecValTok{0}\NormalTok{, n, n)}
    \NormalTok{b[}\KeywordTok{upper.tri} \NormalTok{(b, }\DataTypeTok{diag =} \OtherTok{TRUE}\NormalTok{)] <-}\StringTok{ }\NormalTok{a}
    \NormalTok{b <-}\StringTok{ }\NormalTok{b +}\StringTok{ }\KeywordTok{t}\NormalTok{(b)}
    \KeywordTok{diag} \NormalTok{(b) <-}\StringTok{ }\KeywordTok{diag} \NormalTok{(b) /}\StringTok{ }\DecValTok{2}
    \KeywordTok{return} \NormalTok{(b)}
\NormalTok{\}}
\end{Highlighting}
\end{Shaded}

\subsection{C code}\label{c-code}

\begin{Shaded}
\begin{Highlighting}[]
\DataTypeTok{void} \NormalTok{gsweep_(}\DataTypeTok{double} \NormalTok{*, }\DataTypeTok{double} \NormalTok{*, }\DataTypeTok{int} \NormalTok{*, }\DataTypeTok{int} \NormalTok{*, }\DataTypeTok{int} \NormalTok{*, }\DataTypeTok{int} \NormalTok{*, }\DataTypeTok{double} \NormalTok{*, }\DataTypeTok{int} \NormalTok{*);}

\DataTypeTok{void}
\NormalTok{hsweep(}\DataTypeTok{double} \NormalTok{*s, }\DataTypeTok{double} \NormalTok{*trian, }\DataTypeTok{int} \NormalTok{*n, }\DataTypeTok{int} \NormalTok{*m, }\DataTypeTok{int} \NormalTok{*ind, }\DataTypeTok{int} \NormalTok{*np, }\DataTypeTok{double} \NormalTok{*eps) \{}
    \DataTypeTok{int} \NormalTok{ifault = }\DecValTok{0}\NormalTok{, l = }\DecValTok{0}\NormalTok{;}
    \KeywordTok{for} \NormalTok{(}\DataTypeTok{int} \NormalTok{i = }\DecValTok{0}\NormalTok{; i < *np; i++) \{}
        \NormalTok{gsweep_ (s, trian, ind + i, &l, m, n, eps, &ifault);}
    \NormalTok{\}}
\end{Highlighting}
\end{Shaded}

\section{NEWS}\label{news}

\section*{References}\label{references}
\addcontentsline{toc}{section}{References}

\hypertarget{refs}{}

\end{document}
